% Якоря мест подписей и дат (docs/design/signing-route.md, 15.17.3). Эталон файла пака `common/anchors.tex`:
% пак кладёт его копию байт в байт и подключает `% Якоря мест подписей и дат (docs/design/signing-route.md, 15.17.3). Эталон файла пака `common/anchors.tex`:
% пак кладёт его копию байт в байт и подключает `% Якоря мест подписей и дат (docs/design/signing-route.md, 15.17.3). Эталон файла пака `common/anchors.tex`:
% пак кладёт его копию байт в байт и подключает `% Якоря мест подписей и дат (docs/design/signing-route.md, 15.17.3). Эталон файла пака `common/anchors.tex`:
% пак кладёт его копию байт в байт и подключает `\input{../common/anchors}` после `common/meta`.
%
% \hseSignAnchor[<ширина линии>]{<слот>}       перед линией подписи, на её базовой линии
% \hseDateAnchor[<ширина линии>]{<слот>}{<часть>} перед частью строки даты (day, month_word, year, ...)
%
% На PDF ничего не печатает. При выводе страницы пишет позицию в `<jobname>.hsepos` каталога сборки:
%   hsepos 1 <ширина бумаги> <высота бумаги>
%   sign <слот> - <страница> <x> <y> <ширина>
%   date <слот> <часть> <страница> <x> <y> <ширина>
% числа в sp от левого нижнего угла страницы. Имя и часть: [a-z][a-z0-9_]{0,39}.
\makeatletter
\ExplSyntaxOn
\cs_new_eq:NN \hse@savepos  \tex_savepos:D    % pdfTeX и XeTeX: \pdfsavepos; LuaTeX: \savepos
\cs_new_eq:NN \hse@lastxpos \tex_lastxpos:D
\cs_new_eq:NN \hse@lastypos \tex_lastypos:D
\ExplSyntaxOff
\newwrite\hse@posfile
\newcount\hse@page
\AddToHook{shipout/before}{\global\advance\hse@page\@ne}
\AtBeginDocument{\immediate\openout\hse@posfile=\jobname.hsepos\relax
  \immediate\write\hse@posfile{hsepos 1 \number\paperwidth\space\number\paperheight}}
% #1 sign|date, #2 имя, #3 часть (у подписи «-»), #4 ширина линии
\newcommand\hse@anchor[4]{%
  \leavevmode\hse@savepos
  \edef\hse@tmp{\write\hse@posfile{#1 #2 #3 \noexpand\the\hse@page\space
    \noexpand\the\hse@lastxpos\space\noexpand\the\hse@lastypos\space\number\dimexpr#4\relax}}%
  \hse@tmp}
\newcommand\hseSignAnchor[2][0pt]{\hse@anchor{sign}{#2}{-}{#1}}
\newcommand\hseDateAnchor[3][0pt]{\hse@anchor{date}{#2}{#3}{#1}}
\makeatother
` после `common/meta`.
%
% \hseSignAnchor[<ширина линии>]{<слот>}       перед линией подписи, на её базовой линии
% \hseDateAnchor[<ширина линии>]{<слот>}{<часть>} перед частью строки даты (day, month_word, year, ...)
%
% На PDF ничего не печатает. При выводе страницы пишет позицию в `<jobname>.hsepos` каталога сборки:
%   hsepos 1 <ширина бумаги> <высота бумаги>
%   sign <слот> - <страница> <x> <y> <ширина>
%   date <слот> <часть> <страница> <x> <y> <ширина>
% числа в sp от левого нижнего угла страницы. Имя и часть: [a-z][a-z0-9_]{0,39}.
\makeatletter
\ExplSyntaxOn
\cs_new_eq:NN \hse@savepos  \tex_savepos:D    % pdfTeX и XeTeX: \pdfsavepos; LuaTeX: \savepos
\cs_new_eq:NN \hse@lastxpos \tex_lastxpos:D
\cs_new_eq:NN \hse@lastypos \tex_lastypos:D
\ExplSyntaxOff
\newwrite\hse@posfile
\newcount\hse@page
\AddToHook{shipout/before}{\global\advance\hse@page\@ne}
\AtBeginDocument{\immediate\openout\hse@posfile=\jobname.hsepos\relax
  \immediate\write\hse@posfile{hsepos 1 \number\paperwidth\space\number\paperheight}}
% #1 sign|date, #2 имя, #3 часть (у подписи «-»), #4 ширина линии
\newcommand\hse@anchor[4]{%
  \leavevmode\hse@savepos
  \edef\hse@tmp{\write\hse@posfile{#1 #2 #3 \noexpand\the\hse@page\space
    \noexpand\the\hse@lastxpos\space\noexpand\the\hse@lastypos\space\number\dimexpr#4\relax}}%
  \hse@tmp}
\newcommand\hseSignAnchor[2][0pt]{\hse@anchor{sign}{#2}{-}{#1}}
\newcommand\hseDateAnchor[3][0pt]{\hse@anchor{date}{#2}{#3}{#1}}
\makeatother
` после `common/meta`.
%
% \hseSignAnchor[<ширина линии>]{<слот>}       перед линией подписи, на её базовой линии
% \hseDateAnchor[<ширина линии>]{<слот>}{<часть>} перед частью строки даты (day, month_word, year, ...)
%
% На PDF ничего не печатает. При выводе страницы пишет позицию в `<jobname>.hsepos` каталога сборки:
%   hsepos 1 <ширина бумаги> <высота бумаги>
%   sign <слот> - <страница> <x> <y> <ширина>
%   date <слот> <часть> <страница> <x> <y> <ширина>
% числа в sp от левого нижнего угла страницы. Имя и часть: [a-z][a-z0-9_]{0,39}.
\makeatletter
\ExplSyntaxOn
\cs_new_eq:NN \hse@savepos  \tex_savepos:D    % pdfTeX и XeTeX: \pdfsavepos; LuaTeX: \savepos
\cs_new_eq:NN \hse@lastxpos \tex_lastxpos:D
\cs_new_eq:NN \hse@lastypos \tex_lastypos:D
\ExplSyntaxOff
\newwrite\hse@posfile
\newcount\hse@page
\AddToHook{shipout/before}{\global\advance\hse@page\@ne}
\AtBeginDocument{\immediate\openout\hse@posfile=\jobname.hsepos\relax
  \immediate\write\hse@posfile{hsepos 1 \number\paperwidth\space\number\paperheight}}
% #1 sign|date, #2 имя, #3 часть (у подписи «-»), #4 ширина линии
\newcommand\hse@anchor[4]{%
  \leavevmode\hse@savepos
  \edef\hse@tmp{\write\hse@posfile{#1 #2 #3 \noexpand\the\hse@page\space
    \noexpand\the\hse@lastxpos\space\noexpand\the\hse@lastypos\space\number\dimexpr#4\relax}}%
  \hse@tmp}
\newcommand\hseSignAnchor[2][0pt]{\hse@anchor{sign}{#2}{-}{#1}}
\newcommand\hseDateAnchor[3][0pt]{\hse@anchor{date}{#2}{#3}{#1}}
\makeatother
` после `common/meta`.
%
% \hseSignAnchor[<ширина линии>]{<слот>}       перед линией подписи, на её базовой линии
% \hseDateAnchor[<ширина линии>]{<слот>}{<часть>} перед частью строки даты (day, month_word, year, ...)
%
% На PDF ничего не печатает. При выводе страницы пишет позицию в `<jobname>.hsepos` каталога сборки:
%   hsepos 1 <ширина бумаги> <высота бумаги>
%   sign <слот> - <страница> <x> <y> <ширина>
%   date <слот> <часть> <страница> <x> <y> <ширина>
% числа в sp от левого нижнего угла страницы. Имя и часть: [a-z][a-z0-9_]{0,39}.
\makeatletter
\ExplSyntaxOn
\cs_new_eq:NN \hse@savepos  \tex_savepos:D    % pdfTeX и XeTeX: \pdfsavepos; LuaTeX: \savepos
\cs_new_eq:NN \hse@lastxpos \tex_lastxpos:D
\cs_new_eq:NN \hse@lastypos \tex_lastypos:D
\ExplSyntaxOff
\newwrite\hse@posfile
\newcount\hse@page
\AddToHook{shipout/before}{\global\advance\hse@page\@ne}
\AtBeginDocument{\immediate\openout\hse@posfile=\jobname.hsepos\relax
  \immediate\write\hse@posfile{hsepos 1 \number\paperwidth\space\number\paperheight}}
% #1 sign|date, #2 имя, #3 часть (у подписи «-»), #4 ширина линии
\newcommand\hse@anchor[4]{%
  \leavevmode\hse@savepos
  \edef\hse@tmp{\write\hse@posfile{#1 #2 #3 \noexpand\the\hse@page\space
    \noexpand\the\hse@lastxpos\space\noexpand\the\hse@lastypos\space\number\dimexpr#4\relax}}%
  \hse@tmp}
\newcommand\hseSignAnchor[2][0pt]{\hse@anchor{sign}{#2}{-}{#1}}
\newcommand\hseDateAnchor[3][0pt]{\hse@anchor{date}{#2}{#3}{#1}}
\makeatother
